\documentclass[12pt, leqno]{article}
\usepackage[english,russian]{babel}
\usepackage{fontspec}
\usepackage[utf8]{inputenc}
\usepackage{epigraph}
\usepackage{tabularx}
\usepackage[leqno]{amsmath}
\setmainfont{Times New Roman}
\setlength\parindent{0pt}
\usepackage[colorlinks=true,urlcolor=blue]{hyperref}        
\usepackage{minted} 
\usepackage[xetex,final]{graphicx}
\usepackage{float}%"Плавающие" картинки
\usepackage{wrapfig}%Обтекание фигур (таблиц, картинок и прочег
\usepackage{epstopdf}
 
\begin{document}

\section*{Геометрическая прогрессия. Рекурсия}
\begin{flushleft}
\textbf{Вопросы рекурсии:} Рекурентные формулы и реализация
\end{flushleft}

Геометричекая прогрессия — числовая последовательность вида: 

\begin{equation}
 a,\ ar,\ ar^{2},\ ar^{3},\ ar^{4},\ \ldots
\end{equation}

рекурентная форму имеет вид:


\begin{equation}
a_{n}=a\,r^{n-1}
\end{equation}

Общий член арифметической прогрессии $a_{n}=r\,a_{n-1}$ для  $ n>1 $


Сумма первых $ n $ членов геометрической прогрессии  
\begin{equation}
S_{n}={\begin{cases}\sum \limits _{i=1}^{n}b_{i}={\dfrac {b_{1}-b_{1}q^{n}}{1-q}}={\dfrac {b_{1}\left(1-q^{n}\right)}{1-q}},&amp;{\mbox{if }}q\neq 1\\\\nb_{1},&amp;{\mbox{if }}q=1\end{cases}}
\end{equation}
  



Дальше можно смотреть по ссылке:

\href{https://ru.wikipedia.org/wiki/Геометрическая_прогрессия}{Геометрическая прогрессия}


\clearpage 

\noindent Образец кода на лисп (sbcl):  
 
\begin{minted}[fontsize=\small]{lisp}
;;;; Geometric Progression
(defconstant s 7 "cooef")

;; https://allcalc.ru/node/1001
;; https://en.wikipedia.org/wiki/Geometric_progression


;;; This function recursively build pregression
(defun geometricProgression (summ n)
(
;;Count 'cond'-ition statement"                 
    cond ((<= n 0 ) nil)
         (t (cons ( * summ  s)  (geometricProgression   ( * summ  s ) (- n 1)))) 
    )
)

(defun cli/parameter()
;; Get command line argument" 
    (let (
        (args sb-ext:*posix-argv*))
        (car (cdr args)))
)

;; main print
;; input   ./geometricProgression.lisp  <n = 7>
(defun main()
(terpri)
(time (
       format T "~a  SUMM = ~d" 
       (geometricProgression 1 (- ( parse-integer (cli/parameter)) 1))
       (apply '+ (geometricProgression 1 (- ( parse-integer (cli/parameter)) 1) ))
       ))
)
(main)



\end{minted}  

\clearpage 
Ниже представлен результат для $ n = 7$

\begin{figure}[h]
\center{\includegraphics[width=0.8\textwidth,natwidth=610,natheight=642]{./img/geometric.png}}
\caption{Скриншот вывода нашего скрипта}
\label{ris:image}
\end{figure}

  

\end{document}